%==============================================================================
% infufrgs-example.tex — Complete usage example for the infufrgs class
%
% This file is hereby placed in the public domain.
%
% Compile sequence:
%   pdflatex infufrgs-example
%   bibtex   infufrgs-example
%   pdflatex infufrgs-example
%   pdflatex infufrgs-example
%==============================================================================

% Select program (ppgc, pgmicro, cic, ecp, ppgl) and document type
% (diss/mestrado, tese/doutorado, tc/dipl, ti, rp, espec, pep, prop-tese,
%  plano-doutorado).  Add `english' for English-language documents.
\documentclass[cic,diss,portuguese]{infufrgs}

% If using pdfLaTeX, declare the input encoding (UTF-8 is recommended).
% Not needed with XeLaTeX or LuaLaTeX.
\usepackage[utf8]{inputenc}

% For figure inclusion
\usepackage{graphicx}

% Optionally load math packages
% \usepackage{amsmath}

%------------------------------------------------------------------------------
% Document metadata
%------------------------------------------------------------------------------
\title{Um Exemplo de Monografia do Instituto de Informática da UFRGS}
\translatedtitle{Using \LaTeX\ to Prepare Documents at INF/UFRGS}

% Author: \author{Surname}{First Name(s)}
\author{Machado}{Lucas C.}
% Multiple authors are supported:
%\author{Flaumann}{Frida Gutenberg}

% Advisor (title is optional): \advisor[Title]{Surname}{First Name(s)}
\advisor[Prof.~Dr.]{Scheid}{Eder J.}
% Co-advisor is optional:
%\coadvisor[Prof.~Dr.]{Knuth}{Donald Ervin}

% Date of the defence.  Defaults to the current month and year if omitted.
%\date{maio}{2001}

% Location of the defence.  Defaults to Porto Alegre, RS. 
%\location{Itaquaquecetuba}{SP}

%------------------------------------------------------------------------------
% Keywords — use Title Case; up to 10 keywords per language
%------------------------------------------------------------------------------
\keyword{Ciência da Computação}
\keyword{Formatação de documentos}
\keyword{LaTeX}

\translatedkeyword{Computer Science}
\translatedkeyword{Document formatting}
\translatedkeyword{LaTeX}

%------------------------------------------------------------------------------
% Nominata — override individual fields if the bundled defaults are stale.
% Example:
%   \renewcommand{\nominataReit}{Prof\textsuperscript{a}.~Nome Sobrenome}
%   \renewcommand{\nominataReitname}{Reitora}
%------------------------------------------------------------------------------

%==============================================================================
\begin{document}
%==============================================================================

% Title page + CIP (or nominata page)
\maketitle

%------------------------------------------------------------------------------
% Optional front matter
%------------------------------------------------------------------------------
\begin{dedicatoria}
  À minha família.
\end{dedicatoria}

\begin{agradecimentos}
  Agradeço a todos que contribuíram para a realização deste trabalho.
\end{agradecimentos}

%------------------------------------------------------------------------------
% Abstract in the primary language
%------------------------------------------------------------------------------
\begin{abstract}
Este Trabalho de Conclusão de Curso tem como objetivo avaliar e comparar ferramentas de recuperação de dados utilizadas em Computação Forense, com foco em técnicas de data carving. O estudo será conduzido por meio de uma metodologia experimental que envolverá a seleção de diferentes ferramentas, a criação de imagens e datasets representando distintos cenários de perda de arquivos, a execução controlada das ferramentas e a análise comparativa dos resultados obtidos. Serão considerados diferentes tipos e condições de arquivos, permitindo avaliar a capacidade das ferramentas em recuperar dados e identificar suas limitações em diferentes cenários. Como contribuição, o trabalho busca estabelecer uma metodologia de avaliação que possa ser reproduzida e utilizada em estudos futuros, além de fornecer uma análise comparativa das ferramentas avaliadas. Os resultados serão discutidos considerando aspectos como taxa de recuperação, tipos de arquivos recuperados e limitações identificadas, contribuindo para a compreensão da aplicabilidade das técnicas de data carving em contextos de Computação Forense.
\end{abstract}
%------------------------------------------------------------------------------
% Abstract in the second language (translatedabstract)
%------------------------------------------------------------------------------
% \begin{translatedabstract}
%   This document demonstrates the use of the \textit{infufrgs} class for
%   preparing monographs, dissertations and theses at the Institute of
%   Informatics of UFRGS.  The class provides automatic formatting of the
%   title page, table of contents, lists of figures and tables, ABNT-style
%   references, and other front- and back-matter elements.
% \end{translatedabstract}

%------------------------------------------------------------------------------
% Lists
%------------------------------------------------------------------------------

\tableofcontents

%------------------------------------------------------------------------------
% <Comentado>
%------------------------------------------------------------------------------

% \listoffigures

% \listoftables

% % Optional: list of abbreviations
% \begin{listofabbrv}{UFRGS}
%   \item[ABNT] Associação Brasileira de Normas Técnicas
%   \item[INF]  Instituto de Informática
%   \item[PPGC] Programa de Pós-Graduação em Computação
%   \item[UFRGS] Universidade Federal do Rio Grande do Sul
% \end{listofabbrv}

% Optional: list of symbols
% \begin{listofsymbols}{$\alpha$}
%   \item[$\alpha$] Alpha coefficient
% \end{listofsymbols}
%------------------------------------------------------------------------------
% </Comentado>
%------------------------------------------------------------------------------


%==============================================================================
% Text body
%==============================================================================


% ======================================
% Introdução
% ======================================
\chapter{Introdução}
A Computação Forense é muito importante 
 
\section{Contextualização}
% - Crescimento do uso de dispositivos digitais e do volume de dados armazenados
% - Papel da evidência digital em investigações e processos judiciais
% - Posicionamento da Computação Forense como área de estudo
 
\section{Problema e Motivação}
% - Cenários de perda de dados: exclusão, formatação, corrupção, danos ao sistema de arquivos
% - Recuperação de arquivos quando os metadados não estão mais disponíveis
% - Diversidade de ferramentas com resultados distintos
% - Falta de comparações padronizadas e reprodutíveis
 
\section{Justificativa}
% - Relevância prática para peritos e investigadores
% - Relevância acadêmica (lacuna metodológica)
% - Necessidade de critérios objetivos para a escolha de ferramentas
 
\section{Objetivos}
\subsection{Objetivo geral}
\subsection{Objetivos específicos}
%TODO - Transformar em lista
Objetivos específicos: 
  - seleção das ferramentas,
  - criação de imagens e datasets,
  - execução controlada,
  - análise comparativa e proposta de metodologia
 
\section{Contribuições Esperadas}
%TODO - Transformar em lista
  - Metodologia de avaliação reproduzível
  - Análise comparativa das ferramentas
  - Identificação de limitações por cenário
  - Fornecer 
 
\section{Delimitação do Escopo}
% - Foco em data carving
% - Tipos de arquivos e sistemas de arquivos considerados
% - O que não será abordado (memória, rede, dispositivos móveis, criptografia)
 
\section{Visão Geral da Metodologia}
% - Resumo da abordagem experimental

% TODO - Necessário? 
\section{Organização do Trabalho}
% - Descrição sucinta dos capítulos seguintes
 





% ======================================
% Fundamentação Teórica
% ======================================

\chapter{Fundamentação Teórica}
\label{cap:fundamentacao}

\section{Computação Forense}

\subsection{Evidência Digital}
% - Definição
% - Características
% - Integridade e autenticidade
% - Papel na investigação digital

\subsection{Processo de Investigação Digital}
% - Identificação
% - Aquisição
% - Exame
% - Análise
% - Apresentação
% - Introduzir a recuperação de dados como parte do processo

\section{Armazenamento e Organização de Arquivos}

\subsection{Sistemas de Arquivos}
% - Função de um sistema de arquivos
% - Organização dos arquivos e metadados
% - Estruturas utilizadas para localizar os dados

\subsection{Alocação e Armazenamento de Arquivos}
% - Blocos/clusters
% - Metadados
% - Localização dos dados
% - Relação entre metadados e conteúdo do arquivo

\subsection{Sistemas de Arquivos Considerados}
% - Apresentar somente os sistemas relevantes aos experimentos
% - FAT32/exFAT
% - NTFS
% - Ext4
% - Destacar diferenças relevantes para recuperação

\section{Exclusão e Recuperação de Dados}

\subsection{Exclusão de Arquivos}
% - Exclusão lógica
% - O que ocorre com metadados
% - O que ocorre com os dados armazenados
% - Reutilização dos espaços anteriormente ocupados

\subsection{Formatação e Perda de Dados}
% - Formatação rápida
% - Formatação e alteração das estruturas do sistema de arquivos
% - Diferença entre perda de referências e sobrescrita dos dados

\subsection{Métodos de Recuperação}
% - Recuperação baseada em metadados
% - Recuperação baseada no conteúdo
% - Diferenças entre as abordagens
% - Cenários em que cada abordagem pode ser utilizada

\section{File Carving}

\subsection{Fundamentos do File Carving}
% - Conceito
% - Recuperação sem depender diretamente das estruturas do sistema de arquivos
% - Identificação de arquivos pelo conteúdo

\subsection{Assinaturas de Arquivos}
% - Headers
% - Footers
% - Estruturas internas
% - Tipos de arquivo e assinaturas

\subsection{Reconstrução de Arquivos}
% - Identificação dos limites
% - Extração dos dados
% - Reconstrução do arquivo

\subsection{Fragmentação de Arquivos}
% - Arquivos fragmentados
% - Problema da perda da continuidade dos dados
% - Impacto sobre o carving
% - Técnicas para lidar com fragmentação

\subsection{Limitações do File Carving}
% - Ausência de metadados
% - Nomes de arquivos
% - Estrutura de diretórios
% - Fragmentação
% - Falsos positivos
% - Arquivos parcialmente sobrescritos

\section{Ferramentas de Recuperação de Dados}

\subsection{Ferramentas de File Carving}
% - Foremost
% - Scalpel
% - PhotoRec
% - Outras ferramentas selecionadas
% - Princípios de funcionamento

\subsection{Características das Ferramentas}
% - Sistemas suportados
% - Tipos de arquivos
% - Estratégia de recuperação
% - Configuração
% - Forma de operação

\section{Avaliação de Ferramentas de Recuperação}

\subsection{Métricas de Recuperação}
% - Quantidade de arquivos corretamente recuperados
% - Taxa de recuperação
% - Recuperação por tipo de arquivo

\subsection{Integridade dos Arquivos Recuperados}
% - Validação dos arquivos
% - Hash
% - Comparação com arquivos originais

\subsection{Desempenho}
% - Tempo de execução
% - Consumo de recursos, se aplicável

\subsection{Limitações e Critérios de Comparação}
% - O que caracteriza uma recuperação bem-sucedida
% - Critérios utilizados para comparar as ferramentas







%==============================================================================
% Back matter
%==============================================================================
% <Comentado>
% % Appendices start with \appendix (optional \annexname for annexes)
% \appendix

% \chapter{Exemplo de Apêndice}

% Apêndices contêm material elaborado pelo próprio autor.

% \annex

% \chapter{Exemplo de Anexo}

% Anexos contêm material de terceiros reproduzido para documentação.
% </Comentado>

%------------------------------------------------------------------------------
% References (ABNT author-date style)
%------------------------------------------------------------------------------
\bibliographystyle{abntex2-alf}
\bibliography{infufrgs-example}

\end{document}
